%% claudecode-app-layout.tex — the same four roles, relocated on macOS (figure 2 of 3)
%%
%% One of a COHERENT SET OF THREE rig figures sharing a single visual language
%% (see the "SHARED RIG VISUAL LANGUAGE" block below):
%%   1. vscode-layout.tex          — the four roles in one VS Code window
%%   2. claudecode-app-layout.tex  — the same four roles relocated on macOS   (this file)
%%   3. shell-as-tool-gateway.tex  — the shell as the gateway to tools
%%
%% Purpose: show that the rig's four roles do not disappear off the IDE — they
%% relocate to separate macOS hosts. The layout is deliberately PARALLEL to
%% figure 1 (same colours, same quadrant positions) so a reader maps role-to-role
%% at a glance. The "file being edited" role has no dedicated panel here, so it is
%% drawn dashed/soft to signal that it is implicit, not a fixed pane.
%%
%% Dual-use (standalone preview vs. \input into the book): same guard as figure 1.
\ifdefined\tikzpicture\else
  \documentclass[border=10pt]{standalone}
  \usepackage[T1]{fontenc}
  \usepackage[scaled]{helvet}
  \usepackage[varqu]{inconsolata}
  \usepackage{xcolor}
  \usepackage{tikz}
  \usetikzlibrary{positioning,fit,backgrounds,calc,shapes.geometric,shapes.misc,%
                  arrows.meta,decorations.pathreplacing}
  \begin{document}
  \def\rigEND{\end{document}}
\fi

% ======================================================================
%  SHARED RIG VISUAL LANGUAGE  — keep this block byte-identical in all three
%  rig figures (vscode-layout, claudecode-app-layout, shell-as-tool-gateway).
%  Four roles, one colour + one glyph each, used the same way everywhere:
%     file system  → green   + folder glyph
%     file edited  → amber   + document glyph
%     agent        → violet  + chat glyph
%     shell        → slate   + prompt (>_) glyph   [always the DARK element]
%  \definecolor is safe to repeat; the glyph/legend macros are \providecommand
%  so including several of these figures in one document never clashes.
% ======================================================================
\definecolor{roleFsAcc}{HTML}{2E8B67}   % file system — green accent
\definecolor{roleFsFill}{HTML}{E4F2EB}  % file system — soft fill
\definecolor{roleEdAcc}{HTML}{C6902B}   % file edited — amber accent
\definecolor{roleEdFill}{HTML}{F8EDD6}  % file edited — soft fill
\definecolor{roleAgAcc}{HTML}{7A5EA8}   % agent — violet accent
\definecolor{roleAgFill}{HTML}{ECE6F5}  % agent — soft fill
\definecolor{roleShAcc}{HTML}{37424F}   % shell — slate (the dark element)
\definecolor{roleShFill}{HTML}{DFE4EA}  % shell — soft fill (legends/edges)
\definecolor{roleShInk}{HTML}{D7DCE2}   % shell — light text on slate
\definecolor{roleShCmd}{HTML}{8FC7A6}   % shell — command echo (soft green)
\definecolor{rigInk}{HTML}{2C2C2E}      % primary label text
\definecolor{rigMuted}{HTML}{6E6E76}    % captions / secondary text
\definecolor{rigFrame}{HTML}{BFC2C7}    % window / panel borders
\definecolor{rigSep}{HTML}{D9D9DE}      % thin separators
\definecolor{rigChrome}{HTML}{ECECEE}   % window chrome (title bars)
\definecolor{rigCanvas}{HTML}{FFFFFF}   % editor / page white
%
% ---- role glyphs (tiny silhouettes, ~0.34cm wide, drawn centred at (x,y)) ----
\providecommand{\rigFolder}[3]{% (x,y,colour)
  \begin{scope}[shift={(#1,#2)}]
    \fill[#3] (-0.17,0.02) rectangle (-0.02,0.11);
    \fill[#3] (-0.17,-0.11) rectangle (0.17,0.06);
  \end{scope}}
\providecommand{\rigDoc}[3]{% (x,y,colour)
  \begin{scope}[shift={(#1,#2)}]
    \fill[#3] (-0.12,-0.13) -- (-0.12,0.13) -- (0.04,0.13)
              -- (0.12,0.05) -- (0.12,-0.13) -- cycle;
  \end{scope}}
\providecommand{\rigChat}[3]{% (x,y,colour)
  \begin{scope}[shift={(#1,#2)}]
    \fill[#3, rounded corners=2pt] (-0.15,-0.02) rectangle (0.15,0.13);
    \fill[#3] (-0.10,-0.02) -- (-0.03,-0.02) -- (-0.10,-0.12) -- cycle;
  \end{scope}}
\providecommand{\rigPrompt}[4]{% (x,y,boxcolour,inkcolour)
  \begin{scope}[shift={(#1,#2)}]
    \fill[#3, rounded corners=1.5pt] (-0.17,-0.12) rectangle (0.17,0.13);
    \draw[#4, line width=0.9pt, line cap=round, line join=round]
          (-0.09,0.05) -- (-0.02,-0.005) -- (-0.09,-0.065);
    \draw[#4, line width=0.9pt, line cap=round] (0.01,-0.065) -- (0.10,-0.065);
  \end{scope}}
%
% ---- shared role key (identical on every figure) --------------------------
\providecommand{\rigLegend}[2]{% lower-left of the key sits near (#1,#2)
  \begin{scope}[shift={(#1,#2)}, font=\sffamily]
    \node[anchor=east, text=rigMuted, font=\sffamily\scriptsize] at (-0.15,0) {THE RIG};
    % file system
    \fill[roleFsFill] (0.10,-0.16) rectangle (0.42,0.16);
    \draw[roleFsAcc, line width=0.7pt] (0.10,-0.16) rectangle (0.42,0.16);
    \rigFolder{0.26}{0}{roleFsAcc}
    \node[anchor=west, text=rigInk, font=\sffamily\scriptsize] at (0.52,0) {File system};
    % file edited
    \fill[roleEdFill] (2.55,-0.16) rectangle (2.87,0.16);
    \draw[roleEdAcc, line width=0.7pt] (2.55,-0.16) rectangle (2.87,0.16);
    \rigDoc{2.71}{0}{roleEdAcc}
    \node[anchor=west, text=rigInk, font=\sffamily\scriptsize] at (2.97,0) {File being edited};
    % agent
    \fill[roleAgFill] (5.55,-0.16) rectangle (5.87,0.16);
    \draw[roleAgAcc, line width=0.7pt] (5.55,-0.16) rectangle (5.87,0.16);
    \rigChat{5.71}{-0.01}{roleAgAcc}
    \node[anchor=west, text=rigInk, font=\sffamily\scriptsize] at (5.97,0) {Agent};
    % shell (always the dark swatch)
    \fill[roleShAcc, rounded corners=1pt] (7.35,-0.16) rectangle (7.67,0.16);
    \draw[white, line width=0.9pt, line cap=round, line join=round]
          (7.44,0.05) -- (7.51,-0.005) -- (7.44,-0.065);
    \draw[white, line width=0.9pt, line cap=round] (7.54,-0.065) -- (7.63,-0.065);
    \node[anchor=west, text=rigInk, font=\sffamily\scriptsize] at (7.77,0) {Shell};
  \end{scope}}
% ======================================================================
%  END SHARED BLOCK
% ======================================================================

% ---- macOS chrome (local to this figure; not part of the shared language) ----
\definecolor{macRed}{HTML}{ED6A5E}
\definecolor{macYel}{HTML}{F5BF4F}
\definecolor{macGrn}{HTML}{61C555}
\definecolor{macTitle}{HTML}{F5F6F8}   % window title bar
\definecolor{macMenu}{HTML}{F4F5F7}    % top menu bar
\definecolor{macDeskA}{HTML}{D9E2EE}   % wallpaper (top)
\definecolor{macDeskB}{HTML}{C3D2E4}   % wallpaper (bottom)

\begin{tikzpicture}[
    font=\sffamily,
    x=1cm, y=1cm,
    lbl/.style={anchor=north west, align=left, inner sep=0pt, text=rigInk},
  ]

  % ---- figure super-title ----------------------------------------------------
  \node[anchor=south west, font=\sffamily\bfseries, text=rigInk] at (0,10.18)
    {The same four roles, relocated on macOS — the Claude Code app};
  \node[anchor=south east, font=\sffamily\scriptsize, text=rigMuted] at (16,10.22)
    {same roles, separate apps};

  % ---- desktop wallpaper + soft shadow ---------------------------------------
  \begin{scope}[on background layer]
    \fill[black!9, rounded corners=6pt] (0.10,-0.12) rectangle (16.10,9.88);
  \end{scope}
  \shade[top color=macDeskA, bottom color=macDeskB, rounded corners=6pt] (0,0) rectangle (16,10);
  \draw[rigFrame, line width=0.8pt, rounded corners=6pt] (0,0) rectangle (16,10);

  % ---- macOS menu bar (the "this is macOS" cue) ------------------------------
  \begin{scope}
    \clip[rounded corners=6pt] (0,0) rectangle (16,10);
    \fill[macMenu] (0,9.55) rectangle (16,10);
  \end{scope}
  \draw[rigSep, line width=0.4pt] (0,9.55) -- (16,9.55);
  % apple mark (small silhouette)
  \fill[rigInk] (0.34,9.76) ellipse (0.085 and 0.10);
  \fill[macMenu] (0.45,9.80) circle (0.06);                       % the "bite"
  \fill[rigInk] (0.36,9.90) ellipse (0.045 and 0.03);             % leaf
  \node[anchor=west, font=\sffamily\scriptsize\bfseries, text=rigInk] at (0.52,9.77) {Claude Code};
  \node[anchor=west, font=\sffamily\scriptsize, text=rigMuted] at (2.0,9.77)
    {File\quad Edit\quad View\quad Window\quad Help};
  \node[anchor=east, font=\sffamily\scriptsize, text=rigInk] at (15.8,9.77) {Fri\ \ 3:14\,PM};
  % tiny battery + wifi cues to the left of the clock
  \draw[rigInk, line width=0.5pt, rounded corners=0.5pt] (12.95,9.71) rectangle (13.23,9.83);
  \fill[rigInk] (12.95,9.71) rectangle (13.14,9.83);
  \fill[rigInk] (13.23,9.735) rectangle (13.26,9.805);
  \foreach \r in {0.09,0.15,0.21}{\draw[rigInk, line width=0.5pt] (13.62,9.70) ++(140:\r) arc (140:40:\r);}
  \fill[rigInk] (13.62,9.70) circle (0.02);

  % =====================================================================
  %  FILE SYSTEM — Finder (left column, green) — parallels fig 1 Explorer
  % =====================================================================
  \begin{scope}
    \clip[rounded corners=5pt] (0.4,1.0) rectangle (3.5,9.25);
    \fill[roleFsFill] (0.4,1.0) rectangle (3.5,9.25);           % green tint across the whole pane
    \fill[macTitle]  (0.4,8.75) rectangle (3.5,9.25);            % title bar
  \end{scope}
  \draw[roleFsAcc, line width=1.1pt, rounded corners=5pt] (0.4,1.0) rectangle (3.5,9.25);
  \draw[rigSep, line width=0.4pt] (0.4,8.75) -- (3.5,8.75);
  \draw[rigSep, line width=0.4pt] (0.4,8.2) -- (3.5,8.2);       % caps the File system header band
  \draw[rigSep, line width=0.4pt] (1.5,1.0) -- (1.5,8.2);       % sidebar/files divider, stops below the header
  \fill[macRed] (0.60,9.0) circle (0.052);
  \fill[macYel] (0.77,9.0) circle (0.052);
  \fill[macGrn] (0.94,9.0) circle (0.052);
  \rigFolder{1.24}{9.0}{roleFsAcc}
  \node[anchor=west, font=\sffamily\scriptsize\bfseries, text=roleFsAcc] at (1.42,9.0) {Finder};
  % role label
  \node[lbl, text width=3.0cm] at (0.55,8.55)
    {{\small\bfseries\color{roleFsAcc} File system}};
  % sidebar favourites
  \foreach \i/\t in {0/Recents, 1/Documents, 2/Desktop, 3/Downloads}{
    \node[anchor=west, font=\sffamily\tiny, text=rigMuted] at (0.4,7.6-\i*0.36) {\t};}
  % file list (main pane)
  \foreach \i/\t in {0/{app.py}, 1/{utils.py}, 2/{README.md}, 3/{data.csv}, 4/{figures}}{
    \rigDoc{2.0}{7.62-\i*0.42}{rigMuted}
    \node[anchor=west, font=\ttfamily\tiny, text=rigMuted] at (2.2,7.62-\i*0.42) {\t};}

  % =====================================================================
  %  FILE BEING EDITED — no fixed pane (amber, DASHED/soft) — parallels Editor
  % =====================================================================
  \begin{scope}
    \clip[rounded corners=5pt] (3.95,2.95) rectangle (11.45,9.25);
    \fill[roleEdFill] (3.95,2.95) rectangle (11.45,9.25);
    \fill[macTitle] (3.95,8.75) rectangle (11.45,9.25);
  \end{scope}
  \draw[roleEdAcc, line width=1.1pt, rounded corners=5pt, dashed, dash pattern=on 4pt off 3pt]
        (3.95,2.95) rectangle (11.45,9.25);
  \fill[macRed] (4.15,9.0) circle (0.052);
  \fill[macYel] (4.32,9.0) circle (0.052);
  \fill[macGrn] (4.49,9.0) circle (0.052);
  \rigDoc{4.80}{9.0}{roleEdAcc}
  \node[anchor=west, font=\sffamily\scriptsize\bfseries, text=roleEdAcc] at (4.98,9.0) {Preview / default app};
  \node[lbl, text width=7.2cm] at (4.1,8.5)
    {{\small\bfseries\color{roleEdAcc} File being edited}\ {\scriptsize\color{rigMuted} — files open in app or viewer}};
  % faint document preview page (soft — signals "implicit, not a fixed panel")
  \begin{scope}[opacity=0.55]
    \fill[rigCanvas, rounded corners=2pt] (6.15,3.70) rectangle (9.25,7.90);
    \draw[roleEdAcc!55, line width=0.5pt, rounded corners=2pt] (6.15,3.70) rectangle (9.25,7.90);
    \foreach \i/\w in {0/2.4, 1/2.6, 2/1.7, 3/2.6, 4/2.2, 5/2.6, 6/1.3}{
      \fill[roleEdAcc!22, rounded corners=1pt] (6.45,7.40-\i*0.44) rectangle (6.45+\w,7.40-\i*0.44+0.14);}
  \end{scope}
  \node[anchor=north, font=\sffamily\footnotesize\itshape, text=rigMuted] at (7.7,3.45)
    {opens in Preview, TextEdit, Atom, \dots};

  % =====================================================================
  %  AGENT — inside the Claude Code app window (right, violet) — parallels Claude
  % =====================================================================
  \begin{scope}
    \clip[rounded corners=5pt] (11.9,2.95) rectangle (15.7,9.25);
    \fill[roleAgFill] (11.9,2.95) rectangle (15.7,9.25);
    \fill[macTitle] (11.9,8.75) rectangle (15.7,9.25);
  \end{scope}
  \draw[roleAgAcc, line width=1.1pt, rounded corners=5pt] (11.9,2.95) rectangle (15.7,9.25);
  \draw[rigSep, line width=0.4pt] (11.9,8.75) -- (15.7,8.75);
  \fill[macRed] (12.10,9.0) circle (0.052);
  \fill[macYel] (12.27,9.0) circle (0.052);
  \fill[macGrn] (12.44,9.0) circle (0.052);
  \rigChat{12.74}{9.0}{roleAgAcc}
  \node[anchor=west, font=\sffamily\scriptsize\bfseries, text=roleAgAcc] at (12.92,9.0) {Claude Code};
  \node[lbl, text width=3.7cm] at (12.05,8.5)
    {{\small\bfseries\color{roleAgAcc} Agent}\ {\scriptsize\color{rigMuted}— Claude Code app}};
  \node[lbl, text width=3.55cm] at (12.05,8.08)
    {{\scriptsize\color{rigMuted} lives here}};
  % agent bubble
  \fill[rigCanvas, rounded corners=3pt] (12.1,6.05) rectangle (14.7,7.0);
  \draw[roleAgAcc!45, line width=0.4pt, rounded corners=3pt] (12.1,6.05) rectangle (14.7,7.0);
  \fill[roleAgFill, rounded corners=1pt] (12.3,6.47) rectangle (14.45,6.63);
  \fill[roleAgFill, rounded corners=1pt] (12.3,6.21) rectangle (14.0,6.37);
  \node[font=\sffamily\tiny, text=roleAgAcc, anchor=west] at (12.3,6.86) {agent};
  % user bubble
  \fill[roleAgAcc!16, rounded corners=3pt] (12.9,4.75) rectangle (15.5,5.7);
  \draw[roleAgAcc!38, line width=0.4pt, rounded corners=3pt] (12.9,4.75) rectangle (15.5,5.7);
  \fill[roleAgAcc!38, rounded corners=1pt] (13.1,5.17) rectangle (15.3,5.33);
  \fill[roleAgAcc!38, rounded corners=1pt] (13.1,4.91) rectangle (14.6,5.07);
  \node[font=\sffamily\tiny, text=roleAgAcc, anchor=east] at (15.3,5.56) {you};
  % input box
  \draw[rigFrame, line width=0.5pt, rounded corners=3pt] (12.1,3.25) rectangle (15.5,3.95);
  \node[font=\sffamily\scriptsize, text=rigMuted, anchor=west] at (12.25,3.6) {Ask Claude};
  \fill[roleAgAcc, rounded corners=2pt] (14.9,3.37) rectangle (15.4,3.83);

  % =====================================================================
  %  SHELL — Terminal.app (bottom, slate/dark) — parallels fig 1 Terminal
  % =====================================================================
  \begin{scope}
    \clip[rounded corners=5pt] (3.95,1.0) rectangle (15.7,2.55);
    \fill[roleShAcc] (3.95,1.0) rectangle (15.7,2.55);
    \fill[roleShAcc!85] (3.95,2.15) rectangle (15.7,2.55);       % slightly lighter title bar
  \end{scope}
  \draw[roleShAcc, line width=1.1pt, rounded corners=5pt] (3.95,1.0) rectangle (15.7,2.55);
  \draw[roleShInk!30, line width=0.4pt] (3.95,2.15) -- (15.7,2.15);
  \fill[macRed] (4.15,2.36) circle (0.052);
  \fill[macYel] (4.32,2.36) circle (0.052);
  \fill[macGrn] (4.49,2.36) circle (0.052);
  \rigPrompt{4.80}{2.35}{roleShFill}{roleShAcc}
  \node[anchor=west, font=\sffamily\scriptsize\bfseries, text=white] at (4.98,2.36) {Terminal};
  \node[anchor=west, font=\sffamily\scriptsize, text=roleShInk] at (6.05,2.36)
    {{\bfseries\color{white} Shell}\ — Terminal.app: runs commands; access to tools};
  \node[font=\ttfamily\scriptsize, text=roleShInk, anchor=north west] at (4.15,2.1)
    {\textcolor{roleShCmd}{\$} pandoc notes.md -o notes.pdf};
  \node[font=\ttfamily\scriptsize, text=roleShInk, anchor=north west] at (4.15,1.8)
    {\textcolor{roleShCmd}{\$} python analyze.py};
  \node[font=\ttfamily\scriptsize, text=roleShInk, anchor=north west] at (4.15,1.5)
    {done. \textcolor{white}{\rule[-0.02cm]{0.1cm}{0.2cm}}};

  % ---- macOS Dock (parallels fig 1 status bar position) ----------------------
  \begin{scope}
    \fill[white, opacity=0.55, rounded corners=6pt] (5.3,0.16) rectangle (10.7,0.74);
    \draw[white, opacity=0.7, line width=0.4pt, rounded corners=6pt] (5.3,0.16) rectangle (10.7,0.74);
  \end{scope}
  \foreach \i/\c in {0/roleFsAcc, 1/roleShAcc, 2/roleAgAcc, 3/roleEdAcc, 4/macGrn, 5/macRed}{
    \fill[\c, rounded corners=2.5pt] (5.55+\i*0.85,0.26) rectangle (5.55+\i*0.85+0.5,0.66);}

  % =====================================================================
  %  SHARED ROLE KEY
  % =====================================================================
  \rigLegend{4.4}{-0.72}

\end{tikzpicture}

\ifdefined\rigEND\rigEND\fi
